\section{Choice of parameters and failure of duality}
\label{sec:parameters}

We now choose the parameters in their order of dependence.  The field
size is chosen last, so it can absorb the encoding term from
Proposition~\ref{prop:compression}.

\begin{proof}[Proof of Theorem~\ref{thm:main}]
Fix $a,b\ge1$.  Put
\[
 \theta=\frac1{1000a},
 \qquad s=\lceil1/\theta\rceil,
\]
and choose a positive integer $h$ with
$\log3/\log(h+1)\le\theta$.  Such an integer exists because
$\log(h+1)$ tends to infinity.  Choose a positive integer $r$ so that
\begin{equation}
\label{eq:rank-choice}
 r+h-1>8bh^2s.
\end{equation}
As in Section~\ref{sec:partitions}, define
\[
 D_j=\binom{r+j-1}{j}\quad(0\le j\le h),
 \qquad w=2D_h,
 \qquad u=\left\lceil\frac{hw}{\theta}\right\rceil.
\]
The number
\begin{equation}
\label{eq:fixed-encoding-term}
 C=C_{h,\theta,u}
   =\log u+hs2^u\log\left(1+\frac{s2^u}{\theta}\right)
\end{equation}
is now fixed.  Choose a prime $p$ satisfying
\begin{equation}
\label{eq:prime-choice}
 p>\max\left\{
       h,\ h^2u^{2h},\ \exp\left(\frac{8bC}{D_h}\right)
      \right\}.
\end{equation}
All quantities on the right have already been determined, and primes
exceed every finite bound.

Apply Proposition~\ref{prop:robust-directions} with these parameters,
and form the label maps and graph profiles in
Proposition~\ref{prop:separated-rows}.  The row index set $X$ consists
of the symmetric $h$-linear forms on $\F_p^r$, so
\[
 m=|X|=p^{D_h}>1.
\]
Each label map takes at most $q=p^{D_{h-1}}$ values.  The columns are
indexed by $Y=X\times\mathcal T$, where
\[
 \mathcal T=\{T\subseteq[u]:|[u]\setminus T|\le\theta u\},
\]
and the real row vectors satisfy
\begin{equation}
\label{eq:final-row-separation}
 \|R_x-R_{x'}\|_\infty=1\qquad(x\ne x').
\end{equation}
In particular, $Y$ is a finite nonempty set.  The finite field serves
only to construct the indices and label maps; the matrix entries and
the convex bodies below are real.

By Proposition~\ref{prop:compression}, the positive column class has
a uniform $3\theta$ approximation list of cardinality $Q$ satisfying
\[
 \log Q\le hs\log q+C
         =hsD_{h-1}\log p+C.
\]
Corollary~\ref{cor:signed-compression} gives a uniform $6\theta$
approximation list of at most $Q^2$ functions for the absolutely
convex column hull.  Thus Lemma~\ref{lem:matrix-to-bodies} applies
with $\varepsilon=6\theta$ and $t=\theta$.  It gives admissible
bodies in dimension $n=|Y|$,
\begin{equation}
\label{eq:final-bodies}
 K=3\absconv\{R_x:x\in X\}+\theta B_\infty^Y,
 \qquad L=B_\infty^Y,
\end{equation}
such that
\[
 N(K,L)\ge m,
 \qquad N(L^\circ,38\theta K^\circ)\le Q^2.
\]
Since $38\theta=38/(1000a)<a^{-1}$, increasing the covering body
preserves this cover, and therefore
\begin{equation}
\label{eq:final-dual-bound}
 N(L^\circ,a^{-1}K^\circ)\le Q^2.
\end{equation}

It remains to compare the two entropies.  The exact identity
\[
 \frac{D_{h-1}}{D_h}=\frac{h}{r+h-1}
\]
yields
\begin{equation}
\label{eq:entropy-ratio}
 \frac{\log(Q^2)}{\log m}
 \le \frac{2h^2s}{r+h-1}
      +\frac{2C}{D_h\log p}
 <\frac1{4b}+\frac1{4b}
 =\frac1{2b},
\end{equation}
where the strict inequalities follow respectively from
\eqref{eq:rank-choice} and \eqref{eq:prime-choice}.
Combining \eqref{eq:final-dual-bound} with $N(K,L)\ge m>1$ gives
\[
 b\log N(L^\circ,a^{-1}K^\circ)
 \le b\log(Q^2)
 <\frac12\log m
 <\log N(K,L).
\]
Identifying $\R^Y$ with $\R^n$ makes $L=[-1,1]^n$ and proves the
claimed counterexample for the arbitrary proposed constants $a,b$.
\end{proof}

The same estimate gives an asymptotic form of the conclusion.  Fix
$a\ge1$, and hence $\theta,h,s$.  For each positive integer $r$, set
$D_h,u,C$ as above and choose a prime
\[
 p>\max\left\{h,\ h^2u^{2h},\ \exp\left(\frac{rC}{D_h}\right)
        \right\}.
\]
The resulting bodies $K_r,L_r$ satisfy
\[
 0\le
 \frac{\log N(L_r^\circ,a^{-1}K_r^\circ)}
      {\log N(K_r,L_r)}
 \le \frac{2h^2s}{r+h-1}+\frac2r
 \longrightarrow0.
\]
Although $u$ and $C$ may grow rapidly with $r$, both are fixed before
$p$ is chosen.  No upper bound on the resulting ambient dimension is
required for this conclusion.

\subsection{A consequence for convexified packing}

Theorem~\ref{thm:main} also separates ordinary covering entropy from
convexified-packing entropy. For a nonempty bounded set $T$ and an
origin-symmetric convex body $B$, let $\widehat M(T,B)$ be the largest
length of an ordered sequence $x_1,\ldots,x_m\in T$ such that
\[
 (x_j+\operatorname{int}B)\cap\conv\{x_i:i<j\}=\varnothing
 \qquad(1\le j\le m),
\]
with the first convex hull taken to be empty
\cite[Section~2]{AMSTJ2004}.

\begin{corollary}[Separation from convexified packing]
\label{cor:convexified-separation}
For every $A,C\ge1$, there are an integer $n\ge1$ and an
origin-symmetric convex body $K\subset\R^n$ such that, with
$L=[-1,1]^n$,
\[
 \log N(K,L)>C\log\widehat M(K,L/A).
\]
\end{corollary}

\begin{proof}
An ordered $B$-convexly separated sequence has at most one point in
each translate of $B/4$: two such points have difference in
$B/2\subset\operatorname{int}B$, contradicting the separation condition.
Consequently $\widehat M(T,B)\le N(T,B/4)$, with no boundary ambiguity
for covers by closed translates. The universal convexified duality
bound \cite[Theorem~2]{AMSTJ2004}, applied to $K,L/A$, gives
\[
\begin{split}
 \widehat M(K,L/A)
 &\le\widehat M(A L^\circ,K^\circ/2)^2\\
 &=\widehat M(L^\circ,K^\circ/(2A))^2\\
 &\le N(L^\circ,K^\circ/(8A))^2.
\end{split}
\]
Here the equality uses simultaneous dilation of both arguments.
Choose $K,L$ by Theorem~\ref{thm:main} with $a=8A$ and $b=2C$.
Taking logarithms in the displayed bound proves the corollary.
\end{proof}
